\documentclass[11pt,a4paper]{article}

\usepackage[utf8]{inputenc}
\usepackage[T1]{fontenc}
\usepackage{lmodern}
\usepackage[spanish]{babel}
\usepackage[margin=1in]{geometry}
\usepackage{amsmath,amssymb}
\usepackage{booktabs}
\usepackage{graphicx}
\usepackage{hyperref}
\usepackage{tcolorbox}
\tcbuselibrary{listings}

\newtcblisting{rcode}{
    listing only,
    colback=gray!5,
    colframe=gray!20,
    boxrule=0.4pt,
    sharp corners,
    listing options={
        language=R,
        basicstyle=\ttfamily\footnotesize,
        breaklines=true,
        showstringspaces=false,
        literate=
            {á}{{\'a}}1 {é}{{\'e}}1 {í}{{\'{\i}}}1 {ó}{{\'o}}1 {ú}{{\'u}}1
            {Á}{{\'A}}1 {É}{{\'E}}1 {Í}{{\'I}}1 {Ó}{{\'O}}1 {Ú}{{\'U}}1
            {à}{{\`a}}1 {è}{{\`e}}1 {ì}{{\`\i}}1 {ò}{{\`o}}1 {ù}{{\`u}}1
            {À}{{\`A}}1 {È}{{\`E}}1 {Ì}{{\`I}}1 {Ò}{{\`O}}1 {Ù}{{\`U}}1
            {ñ}{{\~n}}1 {Ñ}{{\~N}}1 {ç}{{\c c}}1 {Ç}{{\c C}}1
            {ü}{{\"u}}1 {Ü}{{\"U}}1
    }
}

\title{Sesión 2 : Inferencia Estadística}
\author{
    Daniel Fuertes\\
    %\texttt{student1@example.com}
    \and
    Oscar Moreno\\
    %\texttt{student2@example.com}
    \and
    Marc Torres\\
    %\texttt{student3@example.com}
}
\date{\today}

\begin{document}

\maketitle

\begin{tcolorbox}[colback=gray!20,colframe=gray!20, sharp corners]
1. Enumere cuales son los elementos clave de una prueba de hipótesis y su significado.
\end{tcolorbox}

La hipótesis alternativa que es la que se quiere probar dada un conocimiento previo o una pregunta sobre el tema, la hipótesis nula que es el complementario de la alternativa, 
el estadístico de la prueba que es el que se estimará a partir de la distribución de referencia que es aquella que le asignaremos al problema dada la muestra.
También el p-value como espacio de probabilidad y el z-value como umbral del estadístico para determinar si nuestra observación es improbable en la hipótesis nula, nos servirá para rechazarla o aceptarla.

\begin{tcolorbox}[colback=gray!20,colframe=gray!20, sharp corners]
2. Un elemento clave en toda prueba de hipótesis es la distribución de referencia del estadístico de la prueba. Diga de que maneras podemos conocer (o aproximar) la distribución de referencia para la prueba de hipótesis de querer ver si una media de una muestra
es igual o no a un cierto valor nominal. Especifique como obtendría o cual sería esta distribución de referencia.
\end{tcolorbox}

Dada una muestra de $n$ observaciones, cogeremos $i$ veces muestras de $m$ observaciones con $m<n$ y calcularemos la media muestral $\overline{X}_{m_i}$ de todas ellas. Representadas en un dotline, veremos una distribución de medias la cual si tomamos con un valor $m$ 
suficientemente grande y con un nombre de repeticiones $i$ también grande, tomará forma simétrica y normal con una función de distribución parecida a $N(\mu,\frac{\sigma}{m})$

\begin{tcolorbox}[colback=gray!20,colframe=gray!20, sharp corners]
3. Un amigo me ofrece un piso por valor de 8MPts. Dejando de lado todos los otros factores intervinientes en la decisión de compra de un piso y solo teniendo en cuenta su precio, puedo considerar que
se trata de una buena ocasión para comprar? Qué suposiciones necesito hacer para resolver el problema?
Utilizar los datos proporcionados.
\end{tcolorbox}

Hay que suponer que la muestra tenemos ha sido obtenida de forma aleatoria, independiente, que con las 221 observaciones ya estimamos la población total de pisos a la venta, y la distribución de referencia es la que resulta de dibujar estos precios en una gráfica.
\\
Sólo hay 5 pisos con valor de 8MPtas. o inferior, entonces tenemos que el $p-value = 5/221 = 0,023$ y en efecto, es un piso barato con 95\% de confianza, deberíamos comprarlo. 


\newpage

\begin{tcolorbox}[colback=gray!20,colframe=gray!20, sharp corners]
4. Sabemos que la media de los precios de los pisos (de 3 habitaciones) en l'Eixample es de $\mu = 16,81$ y su desviación tipo es de $\sigma = 5,91$. Suponiendo que la muestra obtenida
en el ejercicio anterior es aleatoria, ¿ podemos asegurar de que se trata de una muestra de pisos de l'Eixample?
\end{tcolorbox}

\[ x = (7.80\ 8.25\ 12.60\ 13.50\ 15.96\ 16.46\ 31.29)\ , \quad \bar{x} = 15.12\ , \quad n = 7
\]
Supongamos que si que hemos obtenido la muestra del eixample de forma independiente y que la distribución de referencia es normal. En estas condiciones vamos a representar la media observada 
\begin{equation*}
    \overline{x}_n \sim N\Big(\mu, \frac{\sigma^2}{n}\Big) = N(16,810 , 4,99)
\end{equation*}

\begin{figure*}
    %\includegraphics[]{}
\end{figure*}

El p-value da 0,22 que es bastante grande, entonces decimos que obtener una muestra como $x$ de los pisos del eixample es posible.



\begin{tcolorbox}[colback=gray!20,colframe=gray!20, sharp corners]
5. El siguiente año, los precios de una muestra aleatoria de pisos de 3hab. en l'Eixample han sido 
\[
y = (13.57 \ 14.80 \ 22.36 \ 22.70 \ 29.29)
\]
¿ Puedo afirmar que no ha habido cambio de precio entre los dos años?
\end{tcolorbox}

\[\bar{y} = 20,54 \ , \quad n = 5\]

Planteamos el test estadístico

\begin{equation*}
    \begin{cases}
        H_0 : \mu_y - \mu = 0 \\
        H_1 : \mu_y - \mu > 0
    \end{cases}
\end{equation*}
En este caso tenemos que ver la probabilidad de obtener $20,54$ o un valor superior en una muestra de 5 elementos.
\begin{figure*}
    % \includegraphics[]{}
\end{figure*}
El p-value nos da $0.078$ por tanto con una confianza del 95\% no podemos decir que el precio haya subido. Con obtener esta media muestral de una muestra de 7 pisos, entonces sí podriamos afirmarlo.

\newpage

\begin{tcolorbox}[colback=gray!20,colframe=gray!20, sharp corners]
6. Calcular el p valor en la misma prueba del problema anterior, usando el método de permutaciones.
\end{tcolorbox}



\begin{tcolorbox}[colback=gray!20,colframe=gray!20, sharp corners]
7. Sabemos que la probabilidad de compra de un producto en el canal internet es de 0,02. En un mes se han conectado 2300 visitantes, de los cuales 94 han comprado nuestro producto, ¿ puedo pensar
que ha habido un incremento en la probabilidad de compra por internet?
\end{tcolorbox}

Comparación de proporciones, muestral vs proporcion poblacional 

\begin{equation*}
p_{T} = 0,02 \ ; \quad n = 2300 \ ; \quad \hat{p} = \frac{94}{2300} = 0,041 \ ;
\qquad
    \begin{cases}
        H_0 : \hat{p} - p_T \leq 0 \\
        H_1 : \hat{p} - p_T > 0
    \end{cases}
\end{equation*}
Como $p_T$ es una proporción poblacional, la distribución de la observada es 
\begin{equation*}
    \hat{p} \sim N \bigg(p_T, \frac{p_T\cdot(1-p_T)}{n}\bigg) =N\bigg(0.02, \frac{0.02 \cdot (0.98)}{2300}\bigg) = N(0,02 , 8,52 \cdot 10^{-6})
\end{equation*}
\begin{equation*}
    z = \frac{\hat{p}- p_T}{\sqrt{\frac{p_T \cdot (1-p_T)}{n}}} = \frac{0,02}{\sqrt{8,52 \cdot 10 ^{-6}}} = 7,15
\end{equation*}
\begin{rcode}
    > pnorm(z_value, lower.tail = FALSE)
    [1] 4.368528e-13
\end{rcode}

Con este p-valor tan pequeño, tenemos que rechazar la hipótesis nula $H_0$ y concluimos que si ha habido un incremento de la probabilidad.

\newpage

\begin{tcolorbox}[colback=gray!20,colframe=gray!20, sharp corners]
8. Por otro lado, en el mismo mes de la pregunta anterior se ha lanzado una campaña de márqueting directo con un target preseleccionado de 1000 clientes potenciales, con la compra del producto en 56 casos. ¿ Podemos 
afirmar que la tasa de respuesta obtenida en el target preseleecionado es mejor que la obtenida por internet?
\end{tcolorbox}

\begin{equation*}
    p_1 = 0,041 \ ; \quad  n_1 = 2300 \ ; \quad p_2 = 56 / 1000 = 0,056 \ ; \quad n_2 = 1000
\end{equation*}
Suponemos que las muestras de clientes online $x_1$ y la de los preseleccionados $x_2$ son independientes. Planteamos las hipótesis 
\begin{equation*}
    \begin{cases}
        H_0 : p_2 - p_1 \leq 0 \\
        H_1 : p_2 - p_1 > 0
    \end{cases}
\end{equation*}
Como ambas proporciones siguen distribuciones binomiales, su diferencia también:
\begin{equation*}
    \hat{p_2} - \hat{p_1} \sim N \bigg(p_2 - p_1, \frac{p_2 \cdot (1 - p_2)}{n_2} + \frac{p_1 \cdot (1 - p_1)}{n_1}\bigg)
    = N(0,015, 6.99 \cdot 10 ^{-5})
\end{equation*}
Si suponemos que $H_0$ es cierta, tenemos que el $z-value$ es 
\begin{equation*}
    z = \frac{\hat{p}_2 - \hat{p}_1}{\sqrt{\frac{\hat{p}_2 (1 - \hat{p}_2)}{n_2}+ \frac{\hat{p}_1(1-\hat{p}_1)}{n_1}}} \sim N(0,1)
\end{equation*}
\begin{rcode}
    > z12 <- (p2-p1)/ sqrt(v1 + v2)
    > z12
    [1] 1.809633
\end{rcode}
Representamos esta distribución normal estándard con el $z-value$ obtenido
\begin{figure*}
    %\includegraphics*[keyvals]{imagefile}
\end{figure*}
\begin{rcode}
    > pnorm(z12, lower.tail = FALSE)
    [1] 0.03517633
\end{rcode}
Con este p-value podemos decir con 95\% de confianza que las proporción de venta a los preseleccionados es mayor que a los online de forma estadísticamente significativa.
\begin{tcolorbox}[colback=gray!20,colframe=gray!20, sharp corners]
9. Un día me encuentro con un amigo al que hace tiempo que no veía, va acompañado por un hijo varón. Me dice pero que tiene dos hijos, cual es la probabilidad de que su otro hijo sea también varón ?
\end{tcolorbox}

La mitad, ya que el sexo de los hijos es una variable binomial independiente entre si. La probabilidad de ambos varones sería de un cuarto si no hubieses visto que por lo menos uno lo es. 


\begin{tcolorbox}[colback=gray!20,colframe=gray!20, sharp corners]
10. En un concurso de televisión había 3 puertas a escoger, sólo detrás de una de ellas había un premio. El concursante escogía una y luego se revelaba una de las otras dos, una que no tenía el premio.
¿ Cuál es la mejor opción para el concursante, mantenerse en su primera opción o cambiar de puerta?
\end{tcolorbox}

Si es mejor cambiarse, porque la probabilidad de premio en la primera elección es $1/3$ y las otras dos puertas es $2/3$. Cuando se revela que una de ellas no es, la probabilidad de $2/3$ se concentra en la puerta no abierta y no escogida inicialmente.

\end{document}

